\subsection*{CU-05: Activar el segundo factor}
\addcontentsline{toc}{subsection}{CU-05: Activar el segundo factor}
\label{cu-05}

\begin{fichacu}{CU-05}{Activar el segundo factor}
  \crudFila{Responsable}{José Eduardo Prior Hernández}
  \crudFila{Actor principal}{Jugador}
  \crudFila{Actores de sistema}{Servidor de partidas, Proveedor del segundo factor, Servicio de correo, Base de datos}
  \crudFila{Prototipos}{9f, 22d}
  \crudFila{Prioridad}{Must (debe estar en la entrega). Categoría (a) Obligatorio: segundo factor de autenticación (\mbox{CON-08}). El \mbox{CU-01} lo pide al iniciar sesión; este caso lo activa.}
  \crudFila{Objetivo}{Permitir al Jugador con sesión abierta activar el segundo factor de autenticación en su cuenta, comprobando antes con un código de prueba que el canal de entrega funciona, para que desde entonces el \mbox{CU-01} le pida un código además de la contraseña.}
  \crudFila{Disparador}{El Jugador selecciona ``activar'' junto al segundo factor en la \texttt{GUI\_\allowbreak{}AccountSettings}.}
  \textbf{Precondiciones} & \cuCelda{\begin{cureglas}
      \item \textbf{\mbox{PRE-01}} El Jugador tiene una \textit{Sesion} abierta y su token es válido.
      \item \textbf{\mbox{PRE-02}} El \textit{Usuario} tiene \texttt{tipo\_\allowbreak{}cuenta} en \texttt{REGISTRADA} y \texttt{estado\_\allowbreak{}cuenta} en \texttt{ACTIVA}.
      \item \textbf{\mbox{PRE-03}} El \textit{Usuario} tiene \texttt{doble\_\allowbreak{}factor\_\allowbreak{}habilitado} en falso. Si está activado, el Jugador puede desactivarlo con el \mbox{CU-06}.
      \item \textbf{\mbox{PRE-04}} El Servidor de partidas está en ejecución y tiene conexión disponible con la Base de datos.
    \end{cureglas}} \\ \hline
  \textbf{Flujo normal} & \cuCelda{\begin{cuenum}[start=1]
      \item El Jugador selecciona ``activar'' junto al segundo factor en la \texttt{GUI\_\allowbreak{}AccountSettings}, que muestra que está desactivado.
      \item El SJM despliega la \texttt{GUI\_\allowbreak{}EnableSecondFactor}, que explica qué es el segundo factor y solicita la contraseña actual, con las opciones ``Continuar'' y ``Cancelar''. \textit{(Ver \mbox{FA-01})}
      \item El Jugador captura su contraseña y da click en ``Continuar''.
      \item El SJM valida que el campo no esté vacío y envía el mensaje \texttt{ENABLE\_\allowbreak{}SECOND\_\allowbreak{}FACTOR\_\allowbreak{}REQUEST} con la contraseña y el token de sesión. \textit{(Ver \mbox{FA-02}, \mbox{EX-02}, \mbox{EX-03}, \mbox{EX-04})}
      \item El Servidor de partidas verifica que el token corresponda a una \textit{Sesion} abierta. \textit{(Ver \mbox{FA-03})}
      \item El Servidor de partidas deriva el resumen de la contraseña recibida con la \texttt{contrasena\_\allowbreak{}sal} del \textit{Usuario} y lo compara contra \texttt{contrasena\_\allowbreak{}hash} (\mbox{RN-02}). \textit{(Ver \mbox{FA-04})}
      \item El Servidor de partidas verifica que \texttt{doble\_\allowbreak{}factor\_\allowbreak{}habilitado} siga en falso. \textit{(Ver \mbox{FA-05})}
      \item El Servidor de partidas invalida los \textit{CodigoVerificacion} pendientes del \textit{Usuario} con propósito \texttt{SEGUNDO\_\allowbreak{}FACTOR}, genera un código de prueba nuevo con su fecha de expiración y guarda únicamente su resumen criptográfico (\mbox{RN-03}, \mbox{RN-04}). \textit{(Ver \mbox{EX-01})}
    \end{cuenum}} \\
   & \cuCelda{\begin{cuenum}[start=9]
      \item El Servidor de partidas entrega el código al Proveedor del segundo factor. \textit{(Ver \mbox{EX-05})}
      \item El Servidor de partidas responde con \texttt{SECOND\_\allowbreak{}FACTOR\_\allowbreak{}TEST\_\allowbreak{}SENT} y la fecha de expiración del código.
      \item El SJM solicita el código de prueba, mostrando el tiempo restante de vigencia, con las opciones ``Verificar'', ``Reenviar código'' y ``Cancelar''. \textit{(Ver \mbox{FA-01}, \mbox{FA-07})}
      \item El Jugador captura el código y da click en ``Verificar''; el SJM envía el mensaje \texttt{CONFIRM\_\allowbreak{}SECOND\_\allowbreak{}FACTOR\_\allowbreak{}REQUEST} con el código y el token de sesión.
      \item El Servidor de partidas verifica que el código no haya expirado ni se haya usado, y compara su resumen con el almacenado (\mbox{RN-04}). \textit{(Ver \mbox{FA-06})}
      \item El Servidor de partidas, en una sola transacción, marca el \textit{CodigoVerificacion} como usado, escribe \texttt{doble\_\allowbreak{}factor\_\allowbreak{}habilitado} en verdadero en el \textit{Usuario} y registra el cambio en la \textit{BitacoraAcceso} con el resultado \texttt{SEGUNDO\_\allowbreak{}FACTOR\_\allowbreak{}ACTIVADO}. \textit{(Ver \mbox{EX-01})}
      \item El Servidor de partidas solicita al Servicio de correo un aviso al \texttt{correo} del \textit{Usuario} informando de la activación, en su \texttt{idioma\_\allowbreak{}preferido} y fuera de la transacción (\mbox{RN-06}). \textit{(Ver \mbox{EX-05})}
      \item El Servidor de partidas responde con \texttt{SECOND\_\allowbreak{}FACTOR\_\allowbreak{}ENABLED}.
    \end{cuenum}} \\
   & \cuCelda{\begin{cuenum}[start=17]
      \item El SJM muestra la \texttt{GUI\_\allowbreak{}MessageSuccess} indicando que el segundo factor está activado y que se pedirá en el próximo inicio de sesión, y vuelve a la \texttt{GUI\_\allowbreak{}AccountSettings}.
      \item Fin del caso de uso.
    \end{cuenum}} \\ \hline
  \textbf{Flujos alternos} & \cuCelda{\textbf{\mbox{FA-01}: El Jugador cancela}\par
    \begin{cuenum}[start=1]
      \item El Jugador selecciona ``Cancelar''.
      \item Si ya se generó un código de prueba, el SJM envía el mensaje \texttt{CANCEL\_\allowbreak{}SECOND\_\allowbreak{}FACTOR\_\allowbreak{}REQUEST} y el Servidor de partidas invalida el \textit{CodigoVerificacion} pendiente.
      \item El SJM vuelve a la \texttt{GUI\_\allowbreak{}AccountSettings} con el segundo factor desactivado.
      \item Fin del caso de uso.
    \end{cuenum}} \\
   & \cuCelda{\textbf{\mbox{FA-02}: La contraseña está vacía}\par
    \begin{cuenum}[start=1]
      \item El SJM resalta el campo e indica que la contraseña es obligatoria, sin enviar nada al Servidor de partidas.
      \item El flujo continúa en el paso 3 del FN.
    \end{cuenum}} \\
   & \cuCelda{\textbf{\mbox{FA-03}: La sesión ya no es válida}\par
    \begin{cuenum}[start=1]
      \item El Servidor de partidas responde con \texttt{SESSION\_\allowbreak{}INVALID}.
      \item El SJM muestra la \texttt{GUI\_\allowbreak{}MessageWarning} indicando que la sesión terminó y despliega la \texttt{GUI\_\allowbreak{}Login}.
      \item Fin del caso de uso.
    \end{cuenum}} \\
   & \cuCelda{\textbf{\mbox{FA-04}: La contraseña actual es incorrecta}\par
    \begin{cuenum}[start=1]
      \item El Servidor de partidas aplica el mismo tratamiento que el \mbox{FA-04} del \mbox{CU-04}: cuenta el intento para el bloqueo escalonado y lo registra en la \textit{BitacoraAcceso} con el resultado \texttt{CONTRASENA\_\allowbreak{}INCORRECTA} (\mbox{RN-07}).
      \item Si la cuenta quedó bloqueada, el Servidor de partidas cierra todas sus \textit{Sesion} y responde con \texttt{ACCOUNT\_\allowbreak{}LOCKED}; el SJM despliega la \texttt{GUI\_\allowbreak{}Login} con el aviso de bloqueo y termina el caso de uso.
      \item En caso contrario, el Servidor de partidas responde con \texttt{SECOND\_\allowbreak{}FACTOR\_\allowbreak{}FAILED} y el motivo \texttt{CONTRASENA\_\allowbreak{}ACTUAL\_\allowbreak{}INCORRECTA}; el SJM limpia el campo y el flujo continúa en el paso 3 del FN.
    \end{cuenum}} \\
   & \cuCelda{\textbf{\mbox{FA-05}: El segundo factor ya está activado}\par
    \begin{cuenum}[start=1]
      \item El Servidor de partidas encuentra \texttt{doble\_\allowbreak{}factor\_\allowbreak{}habilitado} en verdadero: se activó desde otro dispositivo mientras tanto.
      \item El Servidor de partidas responde con \texttt{SECOND\_\allowbreak{}FACTOR\_\allowbreak{}ALREADY\_\allowbreak{}ENABLED} sin generar ningún código.
      \item El SJM muestra el segundo factor como activado en la \texttt{GUI\_\allowbreak{}AccountSettings}.
      \item Fin del caso de uso.
    \end{cuenum}} \\
   & \cuCelda{\textbf{\mbox{FA-06}: El código de prueba es incorrecto o expiró}\par
    \begin{cuenum}[start=1]
      \item Si el código es incorrecto, el Servidor de partidas incrementa en uno el atributo \texttt{intentos} del \textit{CodigoVerificacion}. Si no alcanzó el máximo del \mbox{RN-04}, responde con \texttt{SECOND\_\allowbreak{}FACTOR\_\allowbreak{}INVALID} y los intentos restantes, que el SJM muestra en la \texttt{GUI\_\allowbreak{}MessageWarning}, y el flujo continúa en el paso 12 del FN.
      \item Si el código expiró o se agotaron los intentos, el Servidor de partidas lo marca como invalidado y responde con \texttt{SECOND\_\allowbreak{}FACTOR\_\allowbreak{}EXPIRED}.
      \item El SJM indica que el código ya no es válido y ofrece ``Reenviar código'', que continúa en el \mbox{FA-07}. El segundo factor sigue desactivado.
    \end{cuenum}} \\
   & \cuCelda{\textbf{\mbox{FA-07}: El Jugador pide reenviar el código}\par
    \begin{cuenum}[start=1]
      \item El Jugador selecciona ``Reenviar código'' y el SJM envía el mensaje \texttt{RESEND\_\allowbreak{}SECOND\_\allowbreak{}FACTOR\_\allowbreak{}TEST\_\allowbreak{}REQUEST}.
      \item El Servidor de partidas verifica que hayan transcurrido al menos sesenta segundos desde la generación del código vigente (\mbox{RN-05}). Si no, responde con los segundos que faltan y el flujo continúa en el paso 12 del FN.
      \item El flujo continúa en el paso 8 del FN.
    \end{cuenum}} \\ \hline
  \textbf{Excepciones} & \cuCelda{\textbf{\mbox{EX-01}: Fallo al escribir en la base de datos}\par
    \begin{cuenum}[start=1]
      \item El Servidor de partidas revierte la transacción, de modo que el segundo factor quede como estaba y no haya un código marcado como usado sin el cambio que confirmaba.
      \item El Servidor de partidas registra la excepción en la bitácora del servidor con un identificador de incidencia, sin incluir la contraseña ni el código recibidos.
      \item El Servidor de partidas responde con \texttt{SERVER\_\allowbreak{}ERROR} y el identificador, que el SJM muestra en la \texttt{GUI\_\allowbreak{}MessageError}.
      \item Fin del caso de uso.
    \end{cuenum}} \\
   & \cuCelda{\textbf{\mbox{EX-02}: Se agota el tiempo de espera de la respuesta}\par
    \begin{cuenum}[start=1]
      \item El SJM detecta que transcurrió el tiempo límite sin recibir un mensaje completo.
      \item El SJM no reenvía la solicitud, vuelve a la \texttt{GUI\_\allowbreak{}AccountSettings} y consulta el estado del segundo factor para mostrar si el cambio se aplicó.
      \item Fin del caso de uso.
    \end{cuenum}} \\
   & \cuCelda{\textbf{\mbox{EX-03}: La conexión se interrumpe durante el intercambio}\par
    \begin{cuenum}[start=1]
      \item El SJM detecta el cierre inesperado del socket antes de recibir respuesta.
      \item El SJM registra el fallo en la bitácora local del cliente y muestra la \texttt{GUI\_\allowbreak{}MessageError} indicando que se perdió la conexión.
      \item Fin del caso de uso.
    \end{cuenum}} \\
   & \cuCelda{\textbf{\mbox{EX-04}: El mensaje recibido está malformado}\par
    \begin{cuenum}[start=1]
      \item El SJM descarta el mensaje, cierra la conexión y registra en la bitácora local el contenido truncado.
      \item El SJM muestra la \texttt{GUI\_\allowbreak{}MessageError} indicando que la respuesta no pudo interpretarse.
      \item Fin del caso de uso.
    \end{cuenum}} \\
   & \cuCelda{\textbf{\mbox{EX-05}: El servicio externo no está disponible}\par
    \begin{cuenum}[start=1]
      \item El Servidor de partidas agota los reintentos previstos hacia el Proveedor del segundo factor o el Servicio de correo.
      \item Si falló la entrega del código de prueba, el Servidor de partidas lo invalida, responde con \texttt{EXTERNAL\_\allowbreak{}SERVICE\_\allowbreak{}UNAVAILABLE} y el segundo factor sigue desactivado; el SJM lo indica en la \texttt{GUI\_\allowbreak{}MessageError} y termina el caso de uso.
      \item Si falló el aviso al correo, el Servidor de partidas conserva el cambio, que ya está confirmado, y registra el fallo del aviso en la bitácora del servidor (\mbox{RN-06}).
    \end{cuenum}} \\ \hline
  \textbf{Postcondiciones} & \cuCelda{\begin{cureglas}
      \item \textbf{\mbox{POST-01}} Si el caso termina por el FN, el \textit{Usuario} tiene \texttt{doble\_\allowbreak{}factor\_\allowbreak{}habilitado} en verdadero y el siguiente inicio de sesión le pide el código (\mbox{FA-10} del \mbox{CU-01}).
      \item \textbf{\mbox{POST-02}} Si el caso termina por cualquier otro flujo alterno o excepción, el segundo factor queda como estaba y ningún código de prueba sigue pendiente.
      \item \textbf{\mbox{POST-03}} Las sesiones abiertas de la cuenta no se cierran en ningún flujo, salvo por el bloqueo del \mbox{FA-04}.
      \item \textbf{\mbox{POST-04}} En todos los casos que llegan al servidor, el intento quedó registrado en la \textit{BitacoraAcceso}.
    \end{cureglas}} \\ \hline
  \textbf{Reglas de negocio} & \cuCelda{\begin{cureglas}
      \item \textbf{\mbox{RN-01}} El segundo factor es opcional para cada cuenta, pero el mecanismo está siempre disponible: así se cumple la restricción que lo exige (\mbox{CON-08}) sin cerrar el acceso a quien no puede administrarlo (\mbox{RN-08} del \mbox{CU-01}, \mbox{Q-03}).
      \item \textbf{\mbox{RN-02}} Activar el segundo factor exige la contraseña actual, aunque haya una sesión abierta.
      \item \textbf{\mbox{RN-03}} Los códigos del segundo factor no se almacenan, transmiten ni registran en claro; en la base de datos se guarda solo su resumen criptográfico (\mbox{RN-02} del \mbox{CU-01}).
      \item \textbf{\mbox{RN-04}} El código de prueba sigue las reglas del código de inicio de sesión: vigencia de cinco minutos, un máximo de tres intentos, un solo uso y uno vigente por cuenta (\mbox{RN-09} del \mbox{CU-01}).
      \item \textbf{\mbox{RN-05}} El código de prueba puede reenviarse a lo sumo una vez cada sesenta segundos, como el del inicio de sesión.
      \item \textbf{\mbox{RN-06}} Toda activación se avisa al correo de la cuenta. El aviso queda fuera de la transacción: que falle no deshace el cambio.
      \item \textbf{\mbox{RN-07}} Una contraseña actual incorrecta cuenta como intento fallido para el bloqueo escalonado (\mbox{RN-07} del \mbox{CU-04}).
      \item \textbf{\mbox{RN-08}} El segundo factor se activa solo tras comprobar un código de prueba. Activarlo sin esa comprobación podría dejar al Jugador fuera de su cuenta si el canal no funciona.
      \item \textbf{\mbox{RN-09}} Activar el segundo factor no cierra las sesiones abiertas: protege los accesos futuros, no los presentes.
      \item \textbf{\mbox{RN-10}} Los invitados no usan este caso: no tienen contraseña ni correo.
    \end{cureglas}} \\ \hline
  \textbf{Observación} & \cuCelda{\textit{Sobre por qué activar el segundo factor es un caso propio.}\par
    En la versión anterior era un flujo alterno del caso de cambiar la contraseña o el correo. Activar el segundo factor es otra meta: no cambia ninguna credencial, añade una comprobación al inicio de sesión, tiene su propio código de prueba y no cierra sesiones. Separarlo deja al \mbox{CU-04} con una sola meta y hace visible dónde se cumple la restricción \mbox{CON-08}. Por la misma razón, desactivarlo es otro caso, el \mbox{CU-06}, que además exige un código del segundo factor.} \\ \hline
  \textbf{Datos que toca} & \cuCelda{\begin{cudatos}
      \item \textbf{\textit{Sesion}:} lee por token \textit{(paso 5)}
      \item \textbf{\textit{Usuario}:} lee \texttt{contrasena\_\allowbreak{}hash}, \texttt{contrasena\_\allowbreak{}sal}, \texttt{doble\_\allowbreak{}factor\_\allowbreak{}habilitado}, \texttt{correo}, \texttt{idioma\_\allowbreak{}preferido} \textit{(pasos 6, 7, 15)}
      \item \textbf{\textit{Usuario}:} escribe \texttt{doble\_\allowbreak{}factor\_\allowbreak{}habilitado}; \texttt{intentos\_\allowbreak{}fallidos} y el bloqueo \textit{(paso 14, \mbox{FA-04})}
      \item \textbf{\textit{CodigoVerificacion}:} invalida los pendientes, crea el de prueba, incrementa \texttt{intentos} y lo marca como usado o invalidado \textit{(pasos 8, 13, 14, \mbox{FA-01}, \mbox{FA-06})}
      \item \textbf{\textit{BitacoraAcceso}:} inserta \textit{(paso 14, \mbox{FA-04})}
    \end{cudatos}} \\ \hline
\end{fichacu}
\clearpage
